% T6 - Synthetic corpus transparency paragraph.
% Target slot: Section 5 (Evaluation) of nids2026_patient_sheets.tex,
% addressing the "Synthetic corpus transparency" bullet (line ~262).
% Length: ~190 words (target band 150-200).

\paragraph{Corpus transparency.}
\label{par:corpus_transparency}
The quantitative evaluation runs on a fully synthetic, fixed bilingual transcript pair: one tripartite teleconsultation in English and one in Spanish, both produced by the project's session simulator.
Each candidate generator then produces five sheets per language at $T{=}0.3$ from this fixed pair (so $n{=}5$ in Section~\ref{sec:evaluation} is generation variance under a fixed input, not transcript variance); transcript-coverage robustness is folded into the upstream-error future work (Section~\ref{sec:conclusions}).
A free-text clinical context written by the author is passed to the cloud-hosted \texttt{gpt-5.1} together with a fixed system prompt that constrains the dialogue to the DOCTOR\,/\,PATIENT\,/\,SPECIALIST roles, enforces a seven-stage flow (introduction, symptoms, initial assessment, specialist connection, collaborative diagnosis, treatment, closing), and targets 15--30 words per turn.
The resulting JSON dialogue is rendered to audio with ElevenLabs TTS (one voice per role) and re-ingested through the same transcription stage applied to real recordings, so synthetic and real audio traverse the same upstream pipeline.
Synthetic data is a deliberate transitional choice while ethics-committee approval for real recordings is in progress.
Crucially, the corpus generator (\texttt{gpt-5.1}, cloud, proprietary) shares no weights with the sheet generator (\texttt{gpt-oss:20b}, local, open-weight) under evaluation, nor with the safety validator (\texttt{qwen3:32b}, local, open-weight; the runner-up of the model-selection protocol), avoiding any direct overlap among the three roles --- corpus, generation, and validation --- on which the evaluation is built.